\section{Dividir y conquistar}
\label{sec:dividir-y-conquistar}

Muchos problemas pueden dividirse en varios subproblemas más pequeños y esos subproblemas siguen siendo muy similares al problema original (por lo que pueden dividirse de nuevo). Tras eso, se pueden combinar las soluciones de los subproblemas para obtener la solución al problema original.

Esa es la idea detrás del paradigma de \concept{dividir y conquistar}, que consiste en:
\begin{enumerate}
    \item \concept{Dividir} el problema en subproblemas más pequeños (usualmente por la mitad)
    \item \concept{Conquistar} (resolver) cada subproblema de forma recursiva (o, si el subproblema es muy pequeño, resolverlo directamente)
    \item \concept{Combinar} las soluciones de los subproblemas, obteniendo la solución al problema original.
\end{enumerate}

Se presentan cuatro algoritmos para concretar esta idea.

\setinputbase{diseno/dividir_conquistar}
\inputlist{merge_sort, subarreglo_max, derivados_merge_sort, notes}

\begin{ejercicios}
    \calentamiento{
        \ejercicio{Longest Common Prefix}{https://leetcode.com/problems/longest-common-prefix/}
    }
    \estandar{
        \ejercicio{Kth Largest Element in an Array}{https://leetcode.com/problems/kth-largest-element-in-an-array/}
        \ejercicio{Search a 2D Matrix II}{https://leetcode.com/problems/search-a-2d-matrix-ii/}
    }
    \duros{
        \ejercicio{Median of Two Sorted Arrays}{https://leetcode.com/problems/median-of-two-sorted-arrays/}
        \ejercicio{The Skyline Problem}{https://leetcode.com/problems/the-skyline-problem/}
    }
\end{ejercicios}
